% GENERATED FILE. Edit canonical lessons, then run npm run generate:books.
% canonical-source-hash: fnv1a64:28eda0c5278a3898
% canonical-lessons: TA-C54-mat, TA-C54-basket, TA-C54-knife, TA-C54-bowl, TA-C54-box

\chapter{Things in the House}
\label{ch:ta-things-in-the-house}
\hlchaptermodality{54}

\begin{chapteropening}
\textbf{By the end of this chapter:} \emph{I can name the mat, basket, knife, bowl and box a Tamil house sets out for a guest.}

Complete TA-C54-box and say all five words in order.
\end{chapteropening}

\section[pāy]{\ta{பாய்} (pāy) — a mat}

\label{lesson:TA-C54-mat}



\begin{quote}
Before the new run of words: what did \emph{pātai} mean, and what did \emph{kirāmam} mean?
\end{quote}

\subsection*{You'll want to know: \ta{பாய்}}
\textbf{\ta{பாய்}} (\emph{pāy}) — \textquotedblleft{}a mat\textquotedblright{}.

The woven mat that is unrolled to sleep on and unrolled again for a guest to sit on. Malayalam says \ml{പായ} (\emph{pāya}) and Telugu \te{చాప} (\emph{cāpa}).

\ta{பாய்} is also 'to leap'. Nothing joins the two but the sound, and what sits around the word settles which is meant — the same one-shape-two-jobs arrangement \emph{tuṇi} already showed you.

\begin{grammarlens}[title={What you've built}]
The first of five things a house puts in front of you.
\end{grammarlens}

\subsection*{Your turn}
Say these aloud:

\begin{itemize}
\raggedright
  \item \emph{pāy}
  \item it once more, slowly
  \item \emph{kirāmam}, then \emph{pāy}, the village and what is unrolled for you in it
\end{itemize}

\begin{tcolorbox}[breakable,colback=teal!4,colframe=teal!35!black,title={Before you move on}]
What does \ta{பாய்} mean? (\textquotedblleft{}a mat\textquotedblright{}.) Say it once more.
\end{tcolorbox}
\section[kūṭai]{\ta{கூடை} (kūṭai) — a basket}

\label{lesson:TA-C54-basket}



\begin{quote}
Before the new one: what did \emph{pāy} mean?
\end{quote}

\subsection*{You'll want to know: \ta{கூடை}}
\textbf{\ta{கூடை}} (\emph{kūṭai}) — \textquotedblleft{}a basket\textquotedblright{}.

Woven, usually of palm leaf or split bamboo. Malayalam's basket is \emph{koṭṭa} and Telugu's is \emph{buṭṭa}; Tamil's is \ta{கூடை}.

It is the word for the basket a vendor sets down in front of you on the pavement, and for the one that goes up onto a head to be carried home.

\begin{grammarlens}[title={What you've built}]
Two: something to sit on, and something to carry in.
\end{grammarlens}

\subsection*{Your turn}
Say these aloud:

\begin{itemize}
\raggedright
  \item \emph{kūṭai}
  \item it once more, slowly
  \item \emph{pāy}, then \emph{kūṭai}, and let the long \ta{ஊ} open the second
\end{itemize}

\begin{tcolorbox}[breakable,colback=teal!4,colframe=teal!35!black,title={Before you move on}]
What does \ta{கூடை} mean? (\textquotedblleft{}a basket\textquotedblright{}.) Say it once more.
\end{tcolorbox}
\section[katti]{\ta{கத்தி} (katti) — a knife}

\label{lesson:TA-C54-knife}



\begin{quote}
Before the new one: what did \emph{kūṭai} mean?
\end{quote}

\subsection*{You'll want to know: \ta{கத்தி}}
\textbf{\ta{கத்தி}} (\emph{katti}) — \textquotedblleft{}a knife\textquotedblright{}.

Three sisters and one word. Malayalam \ml{കത്തി}, Telugu \te{కత్తి}, Tamil \ta{கத்தி}. Not merely close — the same word three times, which happens with things every household has had for as long as there have been households.

The doubled \ta{த} in the middle carries the whole difficulty. Hold it: \emph{ka-t-ti}, with the tongue staying where it is between the two.

\begin{grammarlens}[title={What you've built}]
Three, and this one is shared right across the family.
\end{grammarlens}

\subsection*{Your turn}
Say these aloud:

\begin{itemize}
\raggedright
  \item \emph{katti}
  \item it once more, slowly
  \item \emph{kūṭai}, then \emph{katti}, and hold the doubled \ta{த}
\end{itemize}

\begin{tcolorbox}[breakable,colback=teal!4,colframe=teal!35!black,title={Before you move on}]
What does \ta{கத்தி} mean? (\textquotedblleft{}a knife\textquotedblright{}.) Say it once more.
\end{tcolorbox}
\section[kiṇṇam]{\ta{கிண்ணம்} (kiṇṇam) — a bowl}

\label{lesson:TA-C54-bowl}



\begin{quote}
Before the new one: what did \emph{katti} mean?
\end{quote}

\subsection*{You'll want to know: \ta{கிண்ணம்}}
\textbf{\ta{கிண்ணம்}} (\emph{kiṇṇam}) — \textquotedblleft{}a bowl\textquotedblright{}.

The small round one that holds a helping, rather than the pot the food was cooked in. Telugu's \te{గిన్నె} (\emph{ginne}) is close enough to hear the same word underneath.

Malayalam reaches past both and takes the Sanskrit \ml{പാത്രം} (\emph{pātraṁ}), a vessel of any sort. Tamil has that one as well and keeps it general; \ta{கிண்ணம்} stays small.

\begin{grammarlens}[title={What you've built}]
Four: to sit on, to carry in, to cut with, and to eat from.
\end{grammarlens}

\subsection*{Your turn}
Say these aloud:

\begin{itemize}
\raggedright
  \item \emph{kiṇṇam}
  \item it once more, slowly
  \item \emph{katti}, then \emph{kiṇṇam}, and let the \ta{ண} double this time
\end{itemize}

\begin{tcolorbox}[breakable,colback=teal!4,colframe=teal!35!black,title={Before you move on}]
What does \ta{கிண்ணம்} mean? (\textquotedblleft{}a bowl\textquotedblright{}.) Say it once more.
\end{tcolorbox}
\section[peṭṭi]{\ta{பெட்டி} (peṭṭi) — a box}

\label{lesson:TA-C54-box}



\begin{quote}
Before the new one: what did \emph{kiṇṇam} mean?
\end{quote}

\subsection*{You'll want to know: \ta{பெட்டி}}
\textbf{\ta{பெட்டி}} (\emph{peṭṭi}) — \textquotedblleft{}a box\textquotedblright{}.

Malayalam \ml{പെട്ടി}, Telugu \te{పెట్టె}, Tamil \ta{பெட்டி} — the family agrees again, down to the doubled \ta{ட} in the middle. Two words in one short run where all three sisters say the same thing is worth noticing.

A \ta{பெட்டி} is a box of any size: the tin trunk clothes live in, a matchbox, and the carriage of a train, which Tamil calls a \emph{rayil peṭṭi}.

\begin{grammarlens}[title={What you've built}]
Five things a house sets out: something to sit on, to carry in, to cut with, to eat from, and to keep it all in.
\end{grammarlens}

\subsection*{Your turn}
Say these aloud:

\begin{itemize}
\raggedright
  \item \emph{peṭṭi}
  \item it once more, slowly
  \item all five in order — \emph{pāy}, \emph{kūṭai}, \emph{katti}, \emph{kiṇṇam}, \emph{peṭṭi}
\end{itemize}

\begin{tcolorbox}[breakable,colback=teal!4,colframe=teal!35!black,title={Before you move on}]
What does \ta{பெட்டி} mean? (\textquotedblleft{}a box\textquotedblright{}.) Say it once more.
\end{tcolorbox}
